\documentclass[12pt]{article}

% ── Packages ────────────────────────────────────────────────────────────────
\usepackage[margin=1.2in]{geometry}
\usepackage{amsmath, amssymb}
\usepackage{booktabs}
\usepackage{array}
\usepackage{longtable}
\usepackage{tabularx}
\usepackage{xcolor}
\usepackage{titlesec}
\usepackage{enumitem}
\usepackage{hyperref}
\usepackage{microtype}
\usepackage{parskip}
\usepackage{fancyhdr}
\usepackage{mdframed}
\usepackage{multirow}
\usepackage{colortbl}

% ── Colours ──────────────────────────────────────────────────────────────────
\definecolor{darkblue}{HTML}{1F3864}
\definecolor{midblue}{HTML}{2E75B6}
\definecolor{lightblue}{HTML}{D6E4F0}
\definecolor{altrow}{HTML}{EBF3FB}
\definecolor{mod1}{HTML}{EBF3FB}
\definecolor{mod2}{HTML}{EAF4EA}
\definecolor{mod3}{HTML}{FEF9E7}
\definecolor{mod4}{HTML}{FDEDEC}
\definecolor{midrow}{HTML}{F5F5F5}
\definecolor{examrow}{HTML}{F9EBEA}
\definecolor{presentrow}{HTML}{EBF5EB}

% ── Section formatting ───────────────────────────────────────────────────────
\titleformat{\section}
  {\Large\bfseries\color{darkblue}}{\thesection.}{0.6em}{}
  [\vspace{-4pt}\textcolor{midblue}{\rule{\linewidth}{0.8pt}}\vspace{2pt}]
\titleformat{\subsection}
  {\large\bfseries\color{midblue}}{\thesubsection}{0.6em}{}
\titleformat{\subsubsection}
  {\normalsize\bfseries\color{darkblue}}{\thesubsubsection}{0.6em}{}
\titlespacing{\section}{0pt}{18pt}{6pt}
\titlespacing{\subsection}{0pt}{12pt}{4pt}
\titlespacing{\subsubsection}{0pt}{8pt}{2pt}

% ── Header / footer ──────────────────────────────────────────────────────────
\pagestyle{fancy}
\fancyhf{}
\fancyhead[L]{\small\color{midblue}\textit{Causal Inference for Graduate Students in Statistics}}
\fancyhead[R]{\small\color{midblue}\textit{Iowa State University --- 15-Week Syllabus}}
\fancyfoot[C]{\small\thepage}
\renewcommand{\headrulewidth}{0.4pt}
\renewcommand{\headrule}{\hbox to\headwidth{\color{midblue}\leaders\hrule height \headrulewidth\hfill}}

% ── Lists ────────────────────────────────────────────────────────────────────
\setlist[itemize,1]{leftmargin=1.5em, itemsep=2pt, topsep=4pt,
                    label={\color{midblue}\textbullet}}
\setlist[itemize,2]{leftmargin=2.5em, itemsep=1pt, topsep=2pt,
                    label={\color{midblue}$\circ$}}

% ── Coloured boxes ───────────────────────────────────────────────────────────
\newmdenv[
  backgroundcolor=lightblue, linecolor=midblue, linewidth=1.2pt,
  innerleftmargin=12pt, innerrightmargin=12pt,
  innertopmargin=8pt, innerbottommargin=8pt,
  skipabove=6pt, skipbelow=6pt]{keybox}

\newmdenv[
  backgroundcolor=examrow, linecolor=red!60, linewidth=1.2pt,
  innerleftmargin=12pt, innerrightmargin=12pt,
  innertopmargin=8pt, innerbottommargin=8pt,
  skipabove=6pt, skipbelow=6pt]{exambox}

\newmdenv[
  backgroundcolor=presentrow, linecolor=green!50!black, linewidth=1.2pt,
  innerleftmargin=12pt, innerrightmargin=12pt,
  innertopmargin=8pt, innerbottommargin=8pt,
  skipabove=6pt, skipbelow=6pt]{presentbox}

% ── Column types ─────────────────────────────────────────────────────────────
\newcolumntype{L}[1]{>{\raggedright\arraybackslash}p{#1}}
\newcolumntype{C}[1]{>{\centering\arraybackslash}p{#1}}

% ── Hyperref ─────────────────────────────────────────────────────────────────
\hypersetup{colorlinks=true, linkcolor=midblue,
            urlcolor=midblue, citecolor=midblue}

% ────────────────────────────────────────────────────────────────────────────
\begin{document}
% ────────────────────────────────────────────────────────────────────────────

% ── Title page ───────────────────────────────────────────────────────────────
\thispagestyle{empty}
\begin{center}
  \vspace*{1.2cm}
  {\color{darkblue}\rule{\linewidth}{2pt}}\\[12pt]
  {\Huge\bfseries\color{darkblue}
    Causal Inference for\\[6pt] Graduate Students in Statistics}\\[12pt]
  {\color{midblue}\rule{\linewidth}{1pt}}\\[16pt]
  {\large Department of Statistics \quad\textbullet\quad Iowa State University}\\[8pt]
  {\large Instructor: Jae-kwang Kim}\\[8pt]
  {\normalsize \today}\\[20pt]
  {\color{darkblue}\rule{\linewidth}{2pt}}
\end{center}

\vspace{0.8cm}

\begin{mdframed}[
  backgroundcolor=lightblue, linecolor=midblue, linewidth=1.5pt,
  innerleftmargin=14pt, innerrightmargin=14pt,
  innertopmargin=10pt, innerbottommargin=10pt]
\noindent\textbf{\color{darkblue}Course Overview.}
This course develops causal inference as a two-step discipline:
\emph{identification} (translating a causal quantity into an estimable
functional of the observed data distribution using the do-calculus) and
\emph{estimation} (constructing efficient, robust estimators from a finite
sample using semiparametric theory). The do-operator $\mathrm{do}(T{=}t)$
is the organising principle from week one. A central theme is the distinction
\[
  f\!\bigl(y \mid x,\,\mathrm{do}(T{=}t);\theta\bigr)
  \;\neq\;
  f\!\bigl(y \mid x,\, T{=}t\bigr),
\]
which underlies endogeneity and drives every identification strategy studied
in the course.

\medskip
\noindent\textbf{Structure.} The course runs over 15 weeks at Iowa State
University. Week~8 is a midterm examination; Week~15 is reserved for
student final presentations. Instructional content spans Weeks~1--7 and
9--14, organised into four modules.
\end{mdframed}

\newpage
\tableofcontents
\newpage

% ════════════════════════════════════════════════════════════════════════════
\section{Course Philosophy}
% ════════════════════════════════════════════════════════════════════════════

This course is built around one central message: causal inference requires
two distinct steps that must never be conflated.

\begin{itemize}
  \item \textbf{Identification.}  Translating a causal quantity
    $f(y \mid x, \mathrm{do}(T{=}t); \theta)$ into an estimable function
    of the observed data distribution $f(y,t,x,z)$.  This is the
    do-calculus step.

  \item \textbf{Estimation.}  Constructing efficient, robust estimators
    of the identified functional from a finite sample.  This is the
    statistical step.
\end{itemize}

Most existing courses either focus entirely on step~2 (the econometrics
tradition) or step~1 (the computer-science / Pearl tradition).  This course
develops both, with the connection made explicit throughout.

\medskip
A central theme is the distinction between the \emph{interventional}
conditional density $f(y \mid x, \mathrm{do}(T{=}t); \theta)$ and the
\emph{observational} conditional density $f(y \mid x, T{=}t)$. In the
Gaussian linear model:

\begin{keybox}
\[
  \underbrace{f\!\bigl(y \mid \mathrm{do}(T{=}t)\bigr)}_{\displaystyle
    = N(\beta t,\;\sigma^2)}
  \;\neq\;
  \underbrace{f\!\bigl(y \mid T{=}t\bigr)}_{\displaystyle
    = N\!\Bigl(\bigl(\beta + \tfrac{\rho\sigma}{\tau}\bigr)t,\;
              \sigma^2(1-\rho^2)\Bigr)}.
\]
\end{keybox}

The do-operator makes this distinction \emph{syntactically enforced},
whereas the potential-outcomes framework leaves it implicit.

The 15-week format at Iowa State provides breathing room that a compressed
12-week version cannot: the dense graphical theory (DAGs and $d$-separation)
is separated from the do-calculus rules into two distinct weeks, potential
outcomes receive their own full week after the identification machinery is
in place, and the estimation theory in Module~3 benefits from a
dedicated week on regularisation and calibration that connects the classical
material directly to the instructor's research programme.

% ════════════════════════════════════════════════════════════════════════════
\section{Module 1: Foundations (Weeks 1--4)}
% ════════════════════════════════════════════════════════════════════════════

The first module establishes the conceptual and mathematical infrastructure
on which the rest of the course rests. Each of its four weeks is genuinely
self-contained: the first builds physical intuition for the do-operator, the
second develops the graphical language for reading independence claims, the
third translates that language into identification formulas via the three
rules of do-calculus, and the fourth places the potential-outcomes framework
alongside do-calculus so students see where the two traditions agree and
where they diverge.

% ── Week 1 ───────────────────────────────────────────────────────────────────
\subsection{Week 1: The Interventional/Observational Distinction}

This is the conceptual core of the entire course and deserves a full week.
The central question is: \emph{what does it mean to ask ``what would happen
if we set $T = t$?''} The answer is developed through three lenses.

\begin{itemize}
  \item \textbf{Structural equation model.}
    $Y = f(T, U, \varepsilon)$ and the effect of replacing the equation for
    $T$ with the constant $T = t$ (do-operator as equation surgery).

  \item \textbf{Graphical model.}
    The do-operator as graph surgery --- removing all arrows into $T$ in the
    DAG.

  \item \textbf{Potential outcomes model.}
    $Y(t)$ as a random variable on an extended probability space.
\end{itemize}

The key equation motivating the entire course is derived explicitly:
\begin{keybox}
\[
  f\!\bigl(y \mid x,\, T{=}t\bigr)
  \;=\;
  \int f\!\bigl(y \mid x,\,\mathrm{do}(T{=}t),\, U{=}u\bigr)\;
       p(u \mid x, T{=}t)\; du
  \;\neq\;
  f\!\bigl(y \mid x,\,\mathrm{do}(T{=}t)\bigr).
\]
The gap is exactly the endogeneity bias that causal inference must overcome.
\end{keybox}

The Gaussian linear IV model is introduced as a running example, with
explicit calculations showing the bias term $\rho\sigma/\tau$ separating
the two densities.

\smallskip
\noindent\textit{Reading:} Pearl, Glymour \& Jewell (2016), Ch.~1--4;
Hern\'an \& Robins (2020), Ch.~1--3.

% ── Week 2 ───────────────────────────────────────────────────────────────────
\subsection{Week 2: Directed Acyclic Graphs and \textit{d}-Separation}

In the 15-week format, the graphical theory is given a full week before
the do-calculus rules are introduced. This separation allows students to
become fluent in reading conditional independence from DAGs before they
are asked to manipulate interventional expressions.

\begin{itemize}
  \item \textbf{Directed acyclic graphs.}
    Nodes, directed edges, parents, ancestors, descendants. The acyclicity
    condition and why it is necessary for the Markov factorisation.

  \item \textbf{Global Markov condition.}
    Every variable is independent of its non-descendants given its parents;
    the joint density factorises as
    $p(v_1,\ldots,v_k) = \prod_i p(v_i \mid \mathrm{Pa}(v_i))$.

  \item \textbf{Three path structures.}
    Chain $A \to M \to B$ (blocked by conditioning on $M$);
    fork $A \leftarrow M \rightarrow B$ (blocked by conditioning on $M$);
    collider $A \to C \leftarrow B$ (blocked by default; \emph{opened} by
    conditioning on $C$ or any descendant).

  \item \textbf{\textit{d}-Separation.}
    Formal definition and the Verma--Pearl (1988) theorem: $d$-separation
    implies conditional independence in every distribution Markov with
    respect to the graph.

  \item \textbf{The collider trap.}
    Berkson's bias as an instance of conditioning on a collider. Application
    to the IV DAG: why conditioning on $T$ (the collider in
    $Z \to T \leftarrow U$) opens a back-door path from $Z$ to $Y$.

  \item \textbf{Extended example.}
    Full $d$-separation analysis of the IV DAG
    $\{Z \to T, T \to Y, U \to T, U \to Y\}$ for all interesting
    conditioning sets.
\end{itemize}

\noindent\textit{Reading:} Pearl, Glymour \& Jewell (2016), Ch.~2--3;
Pearl (2000), Ch.~1--2.

% ── Week 3 ───────────────────────────────────────────────────────────────────
\subsection{Week 3: The Do-Calculus and Identification Criteria}

With $d$-separation in hand, students are now ready for the algebraic
machinery that translates graph surgery into identification formulas.

\begin{itemize}
  \item \textbf{Mutilated graphs.}
    $\mathcal{G}_{\overline{X}}$ (delete arrows into $X$) and
    $\mathcal{G}_{\underline{X}}$ (delete arrows out of $X$).

  \item \textbf{Three rules of do-calculus (Pearl, 1995).}
    \begin{itemize}
      \item Rule~1 (insertion/deletion of observations):
        $P(y \mid \mathrm{do}(x), z, w) = P(y \mid \mathrm{do}(x), w)$
        if $(Y \perp\!\!\!\perp Z \mid X,W)_{\mathcal{G}_{\overline{X}}}$.
      \item Rule~2 (action/observation exchange):
        $P(y \mid \mathrm{do}(x),\mathrm{do}(z),w)
         = P(y \mid \mathrm{do}(x),z,w)$
        if $(Y \perp\!\!\!\perp Z \mid X,W)_{\mathcal{G}_{\overline{X}\underline{Z}}}$.
      \item Rule~3 (insertion/deletion of actions):
        $P(y \mid \mathrm{do}(x),\mathrm{do}(z),w)
         = P(y \mid \mathrm{do}(x),w)$
        if $(Y \perp\!\!\!\perp Z \mid X,W)_{\mathcal{G}_{\overline{X}\,\overline{Z(W)}}}$.
    \end{itemize}

  \item \textbf{Back-door criterion.}
    Definition, derivation via Rules~2--3, and the adjustment formula:
    \[
      f\!\bigl(y \mid \mathrm{do}(T{=}t)\bigr)
      = \int f(y \mid t,x)\, p(x)\, dx.
    \]
    Why condition~(1) (no descendant of $T$ in the adjustment set) is
    essential.

  \item \textbf{Front-door criterion.}
    Identification through a mediator $M$ when the confounder is
    unobserved:
    \[
      f\!\bigl(y \mid \mathrm{do}(T{=}t)\bigr)
      = \int P(m \mid t) \int P(y \mid t', m)\, P(t')\, dt'\; dm.
    \]

  \item \textbf{Completeness (Shpitser \& Pearl, 2006).}
    The do-calculus is complete: if $P(y \mid \mathrm{do}(t))$ is not
    identifiable by the three rules, it is genuinely non-identified.
\end{itemize}

\noindent\textit{Reading:} Pearl (2000), Ch.~3--4.

% ── Week 4 ───────────────────────────────────────────────────────────────────
\subsection{Week 4: Potential Outcomes and the Connection to Do-Calculus}

With both the do-calculus and its identification criteria fully established,
Week~4 introduces potential outcomes as a complementary language --- powerful
for defining estimands but, as students can now appreciate, silent on the
identification problem the preceding two weeks have solved.

\begin{itemize}
  \item \textbf{Rubin's potential outcomes.}
    $Y(t)$, SUTVA, and the consistency assumption $Y(t) = Y$ when $T = t$.
    Connection to $f(y \mid \mathrm{do}(T{=}t))$ via SUTVA.

  \item \textbf{Ignorability.}
    $(Y(0), Y(1)) \perp\!\!\!\perp T \mid X$ and its relationship to the
    back-door criterion: ignorability is exactly the graphical condition
    that $X$ blocks all back-door paths.

  \item \textbf{Single World Intervention Graphs (SWIGs).}
    Richardson \& Robins (2014): graphical representation of potential
    outcomes that bridges the two frameworks and makes
    cross-world independence assumptions visible.

  \item \textbf{Where the frameworks agree and diverge.}
    Both define causal estimands clearly. The do-calculus enforces the
    interventional/observational distinction syntactically; potential
    outcomes leave it implicit. For likelihood-based inference, the
    do-calculus is indispensable.

  \item \textbf{Module~1 synthesis.}
    The running Gaussian linear IV example is revisited: students can
    now give a complete account of the endogeneity problem in all three
    languages (SEM, DAG, potential outcomes) and understand why the
    do-operator language is the most precise for the likelihood
    construction studied in Week~11.
\end{itemize}

\noindent\textit{Reading:} Richardson \& Robins (2014);
Imbens \& Rubin (2015), Ch.~1--2.

% ════════════════════════════════════════════════════════════════════════════
\section{Module 2: Identification Theory (Weeks 5--7)}
% ════════════════════════════════════════════════════════════════════════════

% ── Week 5 ───────────────────────────────────────────────────────────────────
\subsection{Week 5: Randomization and Its Surrogates}

\begin{itemize}
  \item \textbf{Randomized experiments.}
    Under randomization $\mathrm{do}(T{=}t)$ is directly observed, so
    $f(y \mid \mathrm{do}(T{=}t)) = f(y \mid T{=}t)$.
    Why this equality holds graphically (no arrows into $T$ in the
    experimental design).

  \item \textbf{Propensity score.}
    $e(X) = P(T{=}1 \mid X)$.  Balancing property:
    $Y(t) \perp\!\!\!\perp T \mid e(X)$.
    Justifies back-door adjustment on $e(X)$ alone.
    Estimation via logistic regression and machine learning.

  \item \textbf{Overlap and positivity.}
    The strict overlap condition $0 < e(X) < 1$; practical consequences
    of near-violations and trimming strategies.

  \item \textbf{Limitation.}
    Back-door adjustment requires no unmeasured confounders --- a strong
    and untestable assumption. This motivates the IV strategy of Week~6.
\end{itemize}

\noindent\textit{Reading:} Hern\'an \& Robins (2020), Ch.~7--9.

% ── Week 6 ───────────────────────────────────────────────────────────────────
\subsection{Week 6: Instrumental Variables --- Identification}

The IV assumptions are stated in do-calculus language:
\begin{itemize}
  \item \textbf{Relevance.}  $Z \not\!\perp\!\!\!\perp T$.
  \item \textbf{Exclusion restriction.}
    $f(y \mid x, \mathrm{do}(T{=}t), z) = f(y \mid x, \mathrm{do}(T{=}t))$.
    This is Rule~1 of the do-calculus, not an observational statement.
  \item \textbf{Exogeneity.}  $Z \perp\!\!\!\perp U \mid X$.
\end{itemize}

\medskip
\noindent\textbf{Case study: the pseudo-density problem.}
The paper proposes the marginal density
\[
  f_m(y \mid x, z;\theta,\phi)
  = \int f\!\bigl(y \mid x, \mathrm{do}(T{=}t);\theta\bigr)\,
         g(t \mid x, z;\phi)\, dt,
\]
treating it as the true conditional density of $Y \mid X, Z$.
The true marginal integrates the \emph{observational} density, and in
the Gaussian linear model the discrepancy is the cross-term
$2\beta\rho\sigma\tau$ in the variance.  As a pseudo-MLE, the correct
asymptotic variance is the sandwich form
\[
  \sqrt{n}\,(\hat\theta - \theta_0) \xrightarrow{d}
  N\!\bigl(0,\; J^{-1}\Sigma(J^{-1})^\top\bigr),
  \qquad
  J = E\!\left[\tfrac{\partial^2 \log f_m}{\partial\theta\,\partial\theta^\top}\right],
  \quad
  \Sigma = \mathrm{Var}\!\left[\tfrac{\partial \log f_m}{\partial\theta}\right].
\]

\noindent\textbf{Three fixes compared:}
\begin{itemize}
  \item \textbf{Control function.}  Condition on the first-stage residual
    $\delta = t - E[T \mid x,z;\hat\phi]$.  Corrected density:
    $f(y \mid x, t, \delta;\theta,\rho)
     = N\!\left(\beta t + \tfrac{\rho\sigma}{\tau}\delta,\;
               \sigma^2(1-\rho^2)\right)$.

  \item \textbf{Exponential tilting.}
    $f_{\mathrm{obs}}(y \mid x,t;\theta,\lambda)
     \propto f(y \mid x,\mathrm{do}(T{=}t);\theta)
             \cdot\exp\!\bigl\{\lambda\, s(y,t,x)\bigr\}$.
    $\lambda = 0$ recovers exogeneity; $\lambda = \rho/(\sigma\tau)$
    recovers the control function.

  \item \textbf{Explicit $U$ model.}  Specify $h(u \mid x;\lambda)$,
    integrate out both $T$ and $U$.  Genuine likelihood; variance
    inflation via the Schur complement
    $\mathcal{I}_{\theta\theta}
     - \mathcal{I}_{\theta\lambda}\mathcal{I}_{\lambda\lambda}^{-1}
       \mathcal{I}_{\lambda\theta}$.
\end{itemize}

\noindent\textit{Reading:} Pearl (2000), Ch.~5;
Angrist, Imbens \& Rubin (1996).

% ── Week 7 ───────────────────────────────────────────────────────────────────
\subsection{Week 7: Beyond IV --- Other Identification Strategies}

\begin{itemize}
  \item \textbf{Regression discontinuity.}
    Local randomization interpretation via do-calculus near the cutoff;
    the sharp and fuzzy RD designs; bandwidth selection.

  \item \textbf{Difference-in-differences.}
    The parallel trends assumption
    $E[Y_i(0) \mid D_i{=}1, t{=}1] - E[Y_i(0) \mid D_i{=}1, t{=}0]
     = E[Y_i(0) \mid D_i{=}0, t{=}1] - E[Y_i(0) \mid D_i{=}0, t{=}0]$
    stated in graphical terms; staggered adoption and recent developments.

  \item \textbf{Mediation analysis.}
    Natural direct and indirect effects via
    $f(y \mid \mathrm{do}(T{=}t, M{=}m))$;
    the cross-world independence assumption and its graphical encoding.

  \item \textbf{Sensitivity analysis.}
    Rosenbaum bounds and $E$-values for violations of the identifying
    assumptions; how to quantify the minimum confounding strength that
    would overturn a conclusion.
\end{itemize}

\noindent\textit{Reading:} Imbens \& Rubin (2015), Ch.~23--24.

% ────────────────────────────────────────────────────────────────────────────
\subsection*{Week 8: Midterm Examination}
\addcontentsline{toc}{subsection}{Week 8: Midterm Examination}
% ────────────────────────────────────────────────────────────────────────────

\begin{exambox}
\noindent\textbf{\color{darkblue}Midterm Examination --- Week 8.}
The midterm covers Modules~1--2 (Weeks~1--7).  It is a closed-book,
in-class examination of 80 minutes.  The examination consists of three
parts:

\begin{enumerate}
  \item \textbf{Graphical reasoning (35\%).}
    Given a DAG, determine $d$-separation relationships, identify valid
    adjustment sets, and classify a proposed estimator as back-door,
    front-door, or neither.

  \item \textbf{Do-calculus derivations (35\%).}
    Apply the three rules to derive identification formulas from first
    principles; prove or disprove specific identification claims.

  \item \textbf{Short-answer conceptual questions (30\%).}
    Explain the pseudo-density problem, the exclusion restriction as a
    do-calculus statement, or the distinction between the sandwich
    variance and the Fisher information inverse.
\end{enumerate}

\noindent\textbf{No new content is introduced this week.}
A review session in the lecture period before the examination summarises
the key results of Modules~1--2.
\end{exambox}

% ════════════════════════════════════════════════════════════════════════════
\section{Module 3: Estimation Theory (Weeks 9--13)}
% ════════════════════════════════════════════════════════════════════════════

This module is where expertise in survey sampling, calibration, and
semiparametric theory becomes directly relevant.  The GEC framework,
Bregman projections, and SREG estimators developed in the instructor's
research programme provide a thread of novel methodology alongside the
classical material.  The 15-week format adds a fifth week in this module
(Week~13) dedicated to regularised estimation and calibration --- material
that would otherwise have to be compressed or omitted.

% ── Week 9 ───────────────────────────────────────────────────────────────────
\subsection{Week 9: From Identification to Estimation --- The EIF}

\begin{itemize}
  \item Once identification is established, the causal effect is a
    functional of the observed data distribution: $\psi = \Psi(P)$.

  \item \textbf{Efficient influence function (EIF).}
    For any regular asymptotically linear estimator $\hat\psi$:
    \[
      \sqrt{n}(\hat\psi - \psi_0) \xrightarrow{d} N(0, V), \qquad
      V \geq E_0\bigl[\varphi(O)^2\bigr],
    \]
    where $\varphi$ is the EIF and $E_0[\varphi^2]$ is the
    semiparametric efficiency bound (Bickel et al., 1998).

  \item Derivation of the EIF for the ATE under back-door identification:
    $\varphi(O) = \mu(1,X) - \mu(0,X) + \frac{T(Y-\mu(1,X))}{e(X)}
                  - \frac{(1-T)(Y-\mu(0,X))}{1-e(X)} - \psi$.

  \item \textbf{Connection to GEC.}
    The EIF is the gradient of the functional in the Hilbert space of
    score functions, analogous to the score in parametric models.
    The generalised entropy calibration estimator achieves the
    semiparametric bound by solving a constrained entropy problem
    whose first-order conditions reproduce the EIF score equation.

  \item Influence function arithmetic: linearity, product rule, and
    the delta method for smooth functionals.
\end{itemize}

\noindent\textit{Reading:} Tsiatis (2006); Bickel et al.\ (1998).

% ── Week 10 ──────────────────────────────────────────────────────────────────
\subsection{Week 10: IPW, Outcome Regression, and Doubly Robust Estimation}

\begin{itemize}
  \item \textbf{IPW estimator.}
    \[
      \hat\tau_{\mathrm{IPW}} = \frac{1}{n}\sum_{i=1}^n \left[
        \frac{T_i Y_i}{\hat e(X_i)} - \frac{(1-T_i)Y_i}{1-\hat e(X_i)}
      \right].
    \]
    Consistency and $\sqrt{n}$-rate under correct propensity model;
    efficiency loss relative to EIF.

  \item \textbf{Outcome regression.}
    \[
      \hat\tau_{\mathrm{OR}} = \frac{1}{n}\sum_{i=1}^n
        \bigl[\hat\mu(1, X_i) - \hat\mu(0, X_i)\bigr].
    \]

  \item \textbf{AIPW / doubly robust estimator.}
    \[
      \hat\tau_{\mathrm{DR}} = \hat\tau_{\mathrm{OR}}
        + \frac{1}{n}\sum_{i=1}^n
          \frac{T_i - \hat e(X_i)}{\hat e(X_i)(1-\hat e(X_i))}
          \bigl(Y_i - \hat\mu(T_i, X_i)\bigr).
    \]
    Doubly robust consistency, efficiency, and the relation to the EIF.

  \item \textbf{Connection to calibration weighting.}
    The GEC framework as a generalisation of IPW incorporating
    covariate balance constraints: the calibration weights solve
    $\min_w D_\phi(w \| q)$ subject to $\sum_i w_i h(X_i) = \bar h$,
    reproducing IPW at the entropic divergence.
\end{itemize}

\noindent\textit{Reading:} Bang \& Robins (2005);
Hern\'an \& Robins (2020), Ch.~13.

% ── Week 11 ──────────────────────────────────────────────────────────────────
\subsection{Week 11: Doubly Robust Estimation in the IV Setting}

\begin{itemize}
  \item \textbf{DR IV estimating equation.}
    \[
      U_{\mathrm{DR}}(\theta) = \frac{1}{n}\sum_{i=1}^n
        \Psi(Y_i, T_i, X_i, Z_i;\theta, \hat\phi, \hat\mu) = 0,
    \]
    where $\Psi$ uses the ratio $T/g(X,Z;\hat\phi)$ as the ``IPW
    weight'' for the instrument.

  \item \textbf{Two distinct robustness mechanisms.}
    \begin{itemize}
      \item \emph{Case~(i), correct $f$:} Fisher score identity forces
        $\mu_0 \equiv 0$; exclusion restriction ensures uniform validity
        over $Z$.
      \item \emph{Case~(ii), correct $g$:} Consistency follows from the
        definition of $\mu_0$ alone, for any $f$.
    \end{itemize}

  \item \textbf{Neyman orthogonality.}
    \[
      E_0\!\left[\frac{\partial \log f_m}{\partial\theta} \cdot
                \frac{\partial \log g}{\partial\phi^\top}\right] = 0
      \quad\text{at }(\theta_0,\phi_0).
    \]
    First-stage estimation error does not inflate the asymptotic
    variance of $\hat\theta$ at the $\sqrt{n}$ rate.
\end{itemize}

\noindent\textit{Reading:} Tchetgen Tchetgen (2017);
Chernozhukov et al.\ (2018).

% ── Week 12 ──────────────────────────────────────────────────────────────────
\subsection{Week 12: Double Machine Learning and Cross-Fitting}

\begin{itemize}
  \item \textbf{Double machine learning (DML).}
    Use ML to estimate nuisance functions $e_0$ and $\mu_0$ at a
    potentially slow rate, while retaining $\sqrt{n}$-consistency for
    $\theta_0$ via Neyman orthogonality.

  \item \textbf{Product bias condition.}
    \[
      \|\hat e - e_0\| \cdot \|\hat\mu - \mu_0\| = o_p(n^{-1/2})
    \]
    is sufficient for $\sqrt{n}$-estimation even when each factor
    converges at $o_p(n^{-1/4})$.

  \item \textbf{Cross-fitting in detail.}
    $K$-fold sample splitting; removing overfitting bias of order
    $(\hat e - e_0)(\hat\mu - \mu_0)$; the role of cross-fitting in
    ensuring oracle rates for the causal parameter.

  \item \textbf{SREG estimator} (Kwon, Yang \& Kim, 2026).
    Sample-split regression for high-dimensional survey data as a
    DML-type estimator: splitting resolves the ``double-use'' bias
    from using the same sample for model fitting and estimation.
\end{itemize}

\noindent\textit{Reading:} Chernozhukov et al.\ (2018);
Kwon, Yang \& Kim (2026).

% ── Week 13 ──────────────────────────────────────────────────────────────────
\subsection{Week 13: Regularisation and Calibration in High Dimensions}

\emph{This week is new relative to the 12-week syllabus} and is made
possible by the additional time available at Iowa State.  It bridges
the DML theory of Week~12 to the calibration and Bregman projection
framework developed in the instructor's research programme.

\begin{itemize}
  \item \textbf{High-dimensional covariate adjustment.}
    When $\dim(X) \gg n$, standard back-door adjustment via OLS or
    logistic regression is infeasible.  Regularised estimators
    (LASSO, ridge, elastic net) introduce bias that must be
    de-biased before causal inference.

  \item \textbf{Regularised calibration.}
    The GEC framework extended to high dimensions: the calibration
    weights solve
    \[
      \min_w D_\phi(w \| q)
      \;\text{ s.t. }\;
      \|A^\top(w - q)\|_\infty \leq \lambda,
    \]
    where $A$ is the covariate matrix and $\lambda$ controls the
    balance tolerance.  The H\"older-conjugate penalty structure
    and its connection to the design-optimal estimator.

  \item \textbf{Distributionally robust calibration (DRC).}
    Extend the GEC budget constraint to a Bregman divergence ball:
    $D_\phi(Q \| P) \leq \rho$.  The budget parameter $\rho$ controls
    the robustness-efficiency trade-off; the minimum feasible $\rho$
    ($\rho_{\min}$) and the three-loop computational algorithm.

  \item \textbf{Connection to DML.}
    The calibration weights correspond to an IPW estimator whose
    propensity function satisfies a regularised balance condition.
    Neyman orthogonality carries over; cross-fitting ensures the
    $\sqrt{n}$ rate for the causal parameter even under high-dimensional
    nuisance estimation.
\end{itemize}

\noindent\textit{Reading:} Kwon \& Kim (2024, GEC Biometrics);
Kim \& Kwon (2025, Bregman calibration preprint).

% ════════════════════════════════════════════════════════════════════════════
\section{Module 4: Special Topics (Weeks 14--15*)}
% ════════════════════════════════════════════════════════════════════════════

\emph{Week~15 is reserved for student final presentations;
see Section~\ref{sec:assessment}.  Instructional content in Module~4
runs in Weeks~14 only, with Week~15 presentations drawing on all four
modules.}

% ── Week 14 ──────────────────────────────────────────────────────────────────
\subsection{Week 14: Likelihood-Based Causal Inference and Frontiers}

This week synthesises the entire course by returning to the question:
how do we embed causal reasoning in a likelihood framework?  The
case-study paper is revisited in full, now understood through all the
machinery developed over the preceding 13 weeks.

\begin{itemize}
  \item \textbf{The pseudo-density problem, revisited.}
    \[
      f_m(y \mid z;\theta,\phi)
      = \int f\!\bigl(y \mid \mathrm{do}(T{=}t);\theta\bigr)\,
             g(t \mid z;\phi)\,dt
      \;\neq\; f_{\mathrm{true}}(y \mid z).
    \]
    Full accounting: the information equality $J = \Sigma$ fails; the
    sandwich variance is correct; the doubly robust IV estimator of
    Week~11 restores efficiency.

  \item \textbf{When is the EM algorithm legitimate?}
    The EM weights $w_i^{(k)}(t) \propto g(t \mid x_i, z_i;\hat\phi)
    \cdot f(y_i \mid x_i, \mathrm{do}(T{=}t);\theta^{(k)})$ are
    pseudo-posteriors.  Under the explicit-$U$ model they become
    genuine Bayesian posteriors.

  \item \textbf{Optimal transport and causal inference.}
    Connections to distributionally robust estimation and the Sinkhorn
    divergence calibration framework; the collapse theorem for standard
    entropic OT in calibration contexts and the pivot to Sinkhorn
    divergence calibration.

  \item \textbf{Causal inference for survey data.}
    Combining design-based and model-based approaches; the role of
    survey weights in the do-calculus identification step.

  \item \textbf{Non-parametric identification.}
    The ID algorithm (Shpitser \& Pearl, 2006): a complete algorithm
    for deciding whether $P(y \mid \mathrm{do}(t))$ is identifiable.

  \item \textbf{Open problem.}
    Likelihood-based causal inference without parametric assumptions
    on $U$, using exponential tilting to bridge the do/observe gap.
\end{itemize}

% ────────────────────────────────────────────────────────────────────────────
\subsection*{Week 15: Final Presentations}
\addcontentsline{toc}{subsection}{Week 15: Final Presentations}
% ────────────────────────────────────────────────────────────────────────────

\begin{presentbox}
\noindent\textbf{\color{darkblue}Final Presentations --- Week 15.}
Each student (or team of two) presents their final project in a 15-minute
talk (12 minutes + 3 minutes for questions).  Presentations are scheduled
across the Week~15 lecture periods.

\medskip
\noindent\textbf{Presentation requirements:}
\begin{itemize}
  \item Clearly state the causal question and the identification strategy
    (do-calculus language required).
  \item Describe the estimation approach and explain why it achieves the
    semiparametric efficiency bound (or why it does not).
  \item Present simulation or empirical results, including honest
    uncertainty quantification.
  \item Identify one open problem or limitation of your approach.
\end{itemize}

\noindent\textbf{Evaluation criteria (equally weighted):}
identification rigour, estimation soundness, clarity of presentation,
and quality of the written report (submitted one week before
presentations).
\end{presentbox}

% ════════════════════════════════════════════════════════════════════════════
\section{Weekly Schedule at a Glance}
% ════════════════════════════════════════════════════════════════════════════

\begingroup
\small
\setlength{\tabcolsep}{5pt}
\renewcommand{\arraystretch}{1.35}

\begin{longtable}{C{0.6cm} L{1.5cm} L{5.5cm} L{5.8cm}}
  \toprule
  \rowcolor{darkblue}
  \color{white}\textbf{Wk} &
  \color{white}\textbf{Module} &
  \color{white}\textbf{Topic} &
  \color{white}\textbf{Key Readings} \\
  \midrule
  \endfirsthead
  \toprule
  \rowcolor{darkblue}
  \color{white}\textbf{Wk} &
  \color{white}\textbf{Module} &
  \color{white}\textbf{Topic} &
  \color{white}\textbf{Key Readings} \\
  \midrule
  \endhead
  \bottomrule
  \endfoot

  % Module 1
  \rowcolor{mod1}
  1 & Module 1 & The Interventional/Observational Distinction
    & \textit{Pearl et al.\ (2016) Ch.~1--4;
              Hern\'an \& Robins (2020) Ch.~1--3} \\
  \rowcolor{mod1}
  2 & Module 1 & DAGs and $d$-Separation
    & \textit{Pearl et al.\ (2016) Ch.~2--3;
              Pearl (2000) Ch.~1--2} \\
  \rowcolor{mod1}
  3 & Module 1 & The Do-Calculus and Identification Criteria
    & \textit{Pearl (2000) Ch.~3--4} \\
  \rowcolor{mod1}
  4 & Module 1 & Potential Outcomes and Connection to Do-Calculus
    & \textit{Richardson \& Robins (2014);
              Imbens \& Rubin (2015) Ch.~1--2} \\
  \midrule

  % Module 2
  \rowcolor{mod2}
  5 & Module 2 & Randomization and Its Surrogates
    & \textit{Hern\'an \& Robins (2020) Ch.~7--9} \\
  \rowcolor{mod2}
  6 & Module 2 & Instrumental Variables --- Identification
    & \textit{Pearl (2000) Ch.~5; Angrist et al.\ (1996)} \\
  \rowcolor{mod2}
  7 & Module 2 & Beyond IV: RD, DiD, Mediation, Sensitivity
    & \textit{Imbens \& Rubin (2015) Ch.~23--24} \\
  \midrule

  % Midterm
  \rowcolor{examrow}
  8 & \textit{Exam} & \textbf{Midterm Examination}
    & \textit{Covers Modules 1--2 (Weeks 1--7)} \\
  \midrule

  % Module 3
  \rowcolor{mod3}
  9 & Module 3 & From Identification to Estimation --- The EIF
    & \textit{Tsiatis (2006); Bickel et al.\ (1998)} \\
  \rowcolor{mod3}
  10 & Module 3 & IPW, Outcome Regression, Doubly Robust Estimation
    & \textit{Bang \& Robins (2005);
              Hern\'an \& Robins (2020) Ch.~13} \\
  \rowcolor{mod3}
  11 & Module 3 & Doubly Robust Estimation in the IV Setting
    & \textit{Tchetgen Tchetgen (2017);
              Chernozhukov et al.\ (2018)} \\
  \rowcolor{mod3}
  12 & Module 3 & Double Machine Learning and Cross-Fitting
    & \textit{Chernozhukov et al.\ (2018);
              Kwon, Yang \& Kim (2026)} \\
  \rowcolor{mod3}
  13 & Module 3 & Regularisation and Calibration in High Dimensions
    & \textit{Kwon \& Kim (2024); Kim \& Kwon (2025)} \\
  \midrule

  % Module 4
  \rowcolor{mod4}
  14 & Module 4 & Likelihood-Based Causal Inference and Frontiers
    & \textit{White (1982); Pearl (2000) Ch.~7;
              Shpitser \& Pearl (2006)} \\
  \midrule

  % Presentations
  \rowcolor{presentrow}
  15 & \textit{Pres.} & \textbf{Final Student Presentations}
    & \textit{Written reports due one week prior} \\
  \bottomrule
\end{longtable}
\endgroup

% ════════════════════════════════════════════════════════════════════════════
\section{Assessment}
\label{sec:assessment}
% ════════════════════════════════════════════════════════════════════════════

\begingroup
\small
\setlength{\tabcolsep}{6pt}
\renewcommand{\arraystretch}{1.5}

\begin{tabularx}{\linewidth}{L{3.4cm} C{1.2cm} X}
  \toprule
  \rowcolor{darkblue}
  \color{white}\textbf{Component} &
  \color{white}\textbf{Weight} &
  \color{white}\textbf{Description} \\
  \midrule
  \rowcolor{altrow}
  Problem Sets (4 total) & 25\% &
    One per module (due end of Weeks~4, 7, 13, 14);
    each requires both identification (do-calculus) and
    estimation components \\
  Midterm Exam (Week 8) & 25\% &
    Closed-book, in-class, 80 minutes;
    covers Modules~1--2 (graphical reasoning,
    do-calculus derivations, short conceptual questions) \\
  \rowcolor{altrow}
  Paper Review & 15\% &
    Written critique of a published paper assessing
    whether the do-operator / observational distinction
    is handled correctly;
    due Week~11 \\
  Final Project \& Presentation & 35\% &
    Written report (25\%) due one week before presentations;
    in-class presentation in Week~15 (10\%);
    must develop a new causal inference tool grounded in
    do-calculus identification and semiparametric estimation \\
  \bottomrule
\end{tabularx}
\endgroup

\medskip
\noindent\textbf{Grading scale.}  A: 90--100, B: 80--89, C: 70--79,
D: 60--69, F: below~60. Plus/minus grades are assigned at the discretion
of the instructor.

\medskip
\noindent\textbf{Problem set schedule.}
\begin{itemize}
  \item Problem Set 1 (Module 1): due end of Week~4.
    Topics: $d$-separation, do-calculus derivations, back-door and
    front-door identification.
  \item Problem Set 2 (Module 2): due end of Week~7.
    Topics: IV identification, pseudo-density problem, sensitivity
    analysis.
  \item Problem Set 3 (Module 3): due end of Week~13.
    Topics: EIF derivation, doubly robust estimation, DML product
    bias condition, calibration weighting.
  \item Problem Set 4 (Module 4): due end of Week~14.
    Topics: sandwich variance, EM validity, exponential tilting,
    an applied causal analysis of a real dataset.
\end{itemize}

\noindent\textbf{Paper review.}  The paper review assignment is modelled
on the analysis developed in the course research notes: students select a
published paper and assess whether the paper correctly distinguishes
$f(y \mid \mathrm{do}(T{=}t))$ from $f(y \mid T{=}t)$, whether the
asymptotic theory is correctly stated for the model used, and what fixes
(if any) the paper requires.  The marginal likelihood IV paper is
available as a model template.

% ════════════════════════════════════════════════════════════════════════════
\section{Potential Outcomes vs.\ Do-Operator}
% ════════════════════════════════════════════════════════════════════════════

\begingroup
\small
\setlength{\tabcolsep}{6pt}
\renewcommand{\arraystretch}{1.4}

\begin{tabularx}{\linewidth}{L{4.5cm} X X}
  \toprule
  \rowcolor{darkblue}
  \color{white}\textbf{Feature} &
  \color{white}\textbf{Potential Outcomes $Y(t)$} &
  \color{white}\textbf{Do-Operator $\mathrm{do}(T{=}t)$} \\
  \midrule
  \rowcolor{altrow}
  Defines causal estimands
    & Clear & Clear \\
  Distinguishes interventional vs.\ observational density
    & Implicit only & Syntactically enforced \\
  \rowcolor{altrow}
  Connects to likelihood construction
    & Silent on this & Via back-door / front-door rules \\
  Graph-based reasoning
    & Not built in (requires SWIGs) & Native \\
  \rowcolor{altrow}
  Standard in econometrics / biostatistics
    & Dominant tradition & Growing rapidly \\
  Handles cross-world independence
    & Implicit (SUTVA) & Explicit via graph structure \\
  \bottomrule
\end{tabularx}
\endgroup

% ════════════════════════════════════════════════════════════════════════════
\section{Recommended Core Texts}
% ════════════════════════════════════════════════════════════════════════════

\begingroup
\small
\setlength{\tabcolsep}{6pt}
\renewcommand{\arraystretch}{1.4}

\begin{tabularx}{\linewidth}{L{5.2cm} X}
  \toprule
  \rowcolor{darkblue}
  \color{white}\textbf{Text} &
  \color{white}\textbf{Role in Course} \\
  \midrule
  \rowcolor{altrow}
  \textit{Pearl, Glymour \& Jewell (2016)}
    & Weeks 1--3: do-calculus foundations; primary entry point \\
  \textit{Hern\'an \& Robins (2020)}
    & Weeks 4--10: bridge from identification to statistical
      estimation; freely available online \\
  \rowcolor{altrow}
  \textit{Pearl (2000)}
    & Weeks 2--3, 6: full identification theory and completeness \\
  \textit{Imbens \& Rubin (2015)}
    & Weeks 4--7: potential outcomes perspective on IV and RCTs \\
  \rowcolor{altrow}
  \textit{Tsiatis (2006)}
    & Weeks 9--11: efficient influence functions and semiparametric
      efficiency bounds \\
  \textit{Chernozhukov et al.\ (2018)}
    & Weeks 11--13: double machine learning and Neyman orthogonality \\
  \bottomrule
\end{tabularx}
\endgroup

% ════════════════════════════════════════════════════════════════════════════
\section{What Makes This Course Distinctive}
% ════════════════════════════════════════════════════════════════════════════

\begin{itemize}
  \item \textbf{Identification first.}
    The do-operator is the organising principle from week one, so students
    develop genuine understanding of why identification and estimation are
    different problems.  The 15-week format reinforces this by giving
    DAGs, $d$-separation, and the do-calculus rules each their own week
    rather than compressing them into one.

  \item \textbf{Likelihood-based inference.}
    The connection to the likelihood framework --- where most statisticians
    work --- is made explicit through the analysis of the pseudo-density
    problem.  The key message is that $f_m \neq f_{\mathrm{true}}$ whenever
    $T$ is endogenous, and that this distinction has concrete consequences
    for asymptotic variance estimation.  No other graduate course in
    statistics treats this connection at this level of rigour.

  \item \textbf{Research thread.}
    The instructor's work on GEC, Bregman projections, distributionally
    robust calibration, SREG estimation, and optimal transport provides a
    thread of novel methodology running through Weeks~9--14 that no
    existing textbook covers.  Week~13 (regularisation and calibration
    in high dimensions) is entirely original to this course and develops
    material from the instructor's active research programme.

  \item \textbf{Midterm as diagnostic.}
    Placing the midterm at Week~8 --- after identification theory is
    complete but before estimation theory begins --- serves as a genuine
    diagnostic: students who cannot correctly apply the do-calculus are
    redirected to the foundational material before the more technically
    demanding estimation weeks begin.
\end{itemize}

% ════════════════════════════════════════════════════════════════════════════
\section*{References}
\addcontentsline{toc}{section}{References}
% ════════════════════════════════════════════════════════════════════════════

\begingroup
\small
\begin{description}[leftmargin=1.8em, labelindent=0pt, itemsep=3pt]
  \item Angrist, J.D., Imbens, G.W.\ \& Rubin, D.B.\ (1996).
    Identification of causal effects using instrumental variables.
    \textit{Journal of the American Statistical Association}, 91, 444--455.

  \item Bang, H.\ \& Robins, J.M.\ (2005).
    Doubly robust estimation in missing data and causal inference models.
    \textit{Biometrics}, 61, 962--973.

  \item Bickel, P.J., Klaassen, C.A.J., Ritov, Y.\ \& Wellner, J.A.\
    (1998).
    \textit{Efficient and Adaptive Estimation for Semiparametric Models}.
    Springer.

  \item Chernozhukov, V., Chetverikov, D., Demirer, M., et al.\ (2018).
    Double/debiased machine learning for treatment and structural
    parameters. \textit{Econometrics Journal}, 21, C1--C68.

  \item Hern\'an, M.A.\ \& Robins, J.M.\ (2020).
    \textit{Causal Inference: What If}.
    Chapman \& Hall/CRC.
    \url{https://www.hsph.harvard.edu/miguel-hernan/causal-inference-book/}.

  \item Imbens, G.W.\ \& Rubin, D.B.\ (2015).
    \textit{Causal Inference for Statistics, Social, and Biomedical
    Sciences}.
    Cambridge University Press.

  \item Kim, J.-K.\ \& Kwon, Y.\ (2025).
    Distributionally robust calibration weighting.
    \textit{Preprint}.

  \item Kwon, Y.\ \& Kim, J.-K.\ (2024).
    Generalised entropy calibration.
    \textit{Biometrics}, 80, ujae007.

  \item Kwon, Y., Yang, S.\ \& Kim, J.-K.\ (2026).
    Sample-split regression for high-dimensional survey data.
    \textit{Biometrika} (forthcoming).

  \item Pearl, J.\ (2000).
    \textit{Causality: Models, Reasoning, and Inference}
    (2nd ed., 2009). Cambridge University Press.

  \item Pearl, J., Glymour, M.\ \& Jewell, N.P.\ (2016).
    \textit{Causal Inference in Statistics: A Primer}. Wiley.

  \item Richardson, T.S.\ \& Robins, J.M.\ (2014).
    ACE bound; single world intervention graphs (SWIGs) and
    identification of causal effects.
    \textit{Technical Report}, University of Washington.

  \item Shpitser, I.\ \& Pearl, J.\ (2006).
    Identification of joint interventional distributions in recursive
    semi-Markovian causal models.
    \textit{AAAI}, 21, 1219--1226.

  \item Tchetgen Tchetgen, E.J.\ (2017).
    A general instrumental variable framework for regression analysis
    with outcome missing not at random.
    \textit{Biometrics}, 73, 1123--1131.

  \item Tsiatis, A.A.\ (2006).
    \textit{Semiparametric Theory and Missing Data}. Springer.

  \item Verma, T.\ \& Pearl, J.\ (1988).
    Causal networks: semantics and expressiveness.
    \textit{Proceedings of UAI}, 352--359.

  \item White, H.\ (1982).
    Maximum likelihood estimation of misspecified models.
    \textit{Econometrica}, 50, 1--25.
\end{description}
\endgroup

\end{document}